\section*{References}
{\small\begin{enumerate}[leftmargin=2.4em,labelsep=0.4em,itemsep=2pt,parsep=0pt,label={\arabic*.}]
\item Martin JA, Hamilton BE, Osterman MJK. Births in the United States, 2023. NCHS Data Brief No. 507. Hyattsville, MD: National Center for Health Statistics; 2024. \url{https://doi.org/10.15620/cdc/158789}
\item Martin JA, Osterman MJK, Driscoll AK. Declines in cigarette smoking during pregnancy in the United States, 2016-2021. NCHS Data Brief No. 458. Hyattsville, MD: National Center for Health Statistics; 2023.
\item Liu B, Xu G, Sun Y, Qiu X, Ryckman KK, Yu Y, et al. Maternal cigarette smoking before and during pregnancy and the risk of preterm birth: a dose-response analysis of 25 million mother-infant pairs. PLoS Med. 2020;17(8):e1003158. \url{https://doi.org/10.1371/journal.pmed.1003158}
\item Stock SJ, Bauld L. Maternal smoking and preterm birth: an unresolved health challenge. PLoS Med. 2020;17(9):e1003386. \url{https://doi.org/10.1371/journal.pmed.1003386}
\item Kondracki AJ. Low birthweight in term singletons mediates the association between maternal smoking intensity exposure status and immediate neonatal intensive care unit admission: the E-value assessment. BMC Pregnancy Childbirth. 2020. \url{https://doi.org/10.1186/s12884-020-02981-1}
\item Bang H, Robins JM. Doubly robust estimation in missing data and causal inference models. Biometrics. 2005;61(4):962-73. \url{https://doi.org/10.1111/j.1541-0420.2005.00377.x}
\item Chernozhukov V, Chetverikov D, Demirer M, Duflo E, Hansen C, Newey W, et al. Double/debiased machine learning for treatment and structural parameters. Econom J. 2018;21(1):C1-68. \url{https://doi.org/10.1111/ectj.12097}
\item VanderWeele TJ, Ding P. Sensitivity analysis in observational research: introducing the E-value. Ann Intern Med. 2017;167(4):268-74. \url{https://doi.org/10.7326/M16-2607}
\item National Center for Health Statistics. User guide to the 2023 natality public use file. Hyattsville, MD: NCHS; 2024. \url{https://ftp.cdc.gov/pub/Health_Statistics/NCHS/Dataset_Documentation/DVS/natality/UserGuide2023.pdf}
\item Martin JA, Osterman MJK, Kirmeyer SE, Gregory ECW. Measuring gestational age in vital statistics data: transitioning to the obstetric estimate. Natl Vital Stat Rep. 2015;64(5).
\item Tennant PWG, Murray EJ, Arnold KF, Berrie L, Fox MP, Gadd SC, et al. Use of directed acyclic graphs (DAGs) to identify confounders in applied health research: review and recommendations. Int J Epidemiol. 2021;50(2):620-32. \url{https://doi.org/10.1093/ije/dyaa213}
\item Pearl J. Causal diagrams for empirical research. Biometrika. 1995;82(4):669-88. \url{https://doi.org/10.1093/biomet/82.4.669}
\item Shin D, Song WO. Influence of the adequacy of the Prenatal Care Utilization Index on small-for-gestational-age infants and preterm births in the United States. J Clin Med. 2019;8(6):838. \url{https://doi.org/10.3390/jcm8060838}
\item Joyce T, Racine A, Yunzal-Butler C. Reassessing the WIC effect: evidence from the Pregnancy Nutrition Surveillance System. J Policy Anal Manage. 2008;27(2):277-303.
\item Heaman MI, Newburn-Cook CV, Green CG, Elliott LJ, Helewa ME. Inadequate prenatal care and its association with adverse pregnancy outcomes: a comparison of indices. BMC Pregnancy Childbirth. 2008;8:15.
\item Rubin DB. Estimating causal effects of treatments in randomized and nonrandomized studies. J Educ Psychol. 1974;66(5):688-701.
\item Rosenbaum PR, Rubin DB. The central role of the propensity score in observational studies for causal effects. Biometrika. 1983;70(1):41-55. \url{https://doi.org/10.1093/biomet/70.1.41}
\item Hern\'an MA, Robins JM. Causal inference: what if. Boca Raton: Chapman \& Hall/CRC; 2020.
\item Cole SR, Frangakis CE. The consistency statement in causal inference: a definition or an assumption? Epidemiology. 2009;20(1):3-5. \url{https://doi.org/10.1097/EDE.0b013e31818ef366}
\item Petersen ML, Porter KE, Gruber S, Wang Y, van der Laan MJ. Diagnosing and responding to violations in the positivity assumption. Stat Methods Med Res. 2012;21(1):31-54. \url{https://doi.org/10.1177/0962280210386207}
\item Funk MJ, Westreich D, Wiesen C, St\"urmer T, Brookhart MA, Davidian M. Doubly robust estimation of causal effects. Am J Epidemiol. 2011;173(7):761-7. \url{https://doi.org/10.1093/aje/kwq439}
\item Li F, Morgan KL, Zaslavsky AM. Balancing covariates via propensity score weighting. J Am Stat Assoc. 2018;113(521):390-400. \url{https://doi.org/10.1080/01621459.2016.1260466}
\item Zivich PN, Breskin A. Machine learning for causal inference: on the use of cross-fit estimators. Epidemiology. 2021;32(3):393-401. \url{https://doi.org/10.1097/EDE.0000000000001332}
\item Austin PC, Stuart EA. Moving towards best practice when using inverse probability of treatment weighting (IPTW) using the propensity score to estimate causal treatment effects in observational studies. Stat Med. 2015;34(28):3661-79. \url{https://doi.org/10.1002/sim.6607}
\item Dietz PM, Homa D, England LJ, Burley K, Tong VT, Dube SR, et al. Estimates of nondisclosure of cigarette smoking among pregnant and nonpregnant women of reproductive age in the United States. Am J Epidemiol. 2011;173(3):355-9. \url{https://doi.org/10.1093/aje/kwq381}
\item Howland RE, Mulready-Ward C, Madsen AM, Sackoff J, Nyland-Funke M, Bombard JM, et al. Reliability of reported maternal smoking: comparing the birth certificate to maternal worksheets and prenatal and hospital medical records, New York City and Vermont, 2009. Matern Child Health J. 2015;19(9):1916-24. \url{https://doi.org/10.1007/s10995-015-1722-1}
\item Searles Nielsen S, Dills RL, Glass M, Mueller BA. Accuracy of prenatal smoking data from Washington State birth certificates in a population-based sample with cotinine measurements. Ann Epidemiol. 2014;24(3). \url{https://doi.org/10.1016/j.annepidem.2013.12.008}
\item Jurek AM, Greenland S, Maldonado G, Church TR. Proper interpretation of non-differential misclassification effects: expectations vs observations. Int J Epidemiol. 2005;34(3):680-7. \url{https://doi.org/10.1093/ije/dyi060}
\item Leung M, Kioumourtzoglou MA, Raz R, Weisskopf MG. Bias due to selection on live births in studies of environmental exposures during pregnancy: a simulation study. Environ Health Perspect. 2021;129(4):047001. \url{https://doi.org/10.1289/EHP7961}
\item Patra J, Bakker R, Irving H, Jaddoe VWV, Malini S, Rehm J. Dose-response relationship between alcohol consumption before and during pregnancy and the risks of low birthweight, preterm birth and small for gestational age (SGA): a systematic review and meta-analyses. BJOG. 2011;118(12):1411-21. \url{https://doi.org/10.1111/j.1471-0528.2011.03050.x}
\item Regan AK, Bombard JM, O'Hegarty MM, Smith RA, Tong VT. Adverse birth outcomes associated with prepregnancy and prenatal electronic cigarette use. Obstet Gynecol. 2021;138(1):85-94. \url{https://doi.org/10.1097/AOG.0000000000004432}
\end{enumerate}}
